\begin{table}[!htbp]\centering\scriptsize
\caption{Pre-policy minority-equity adjustment for the initially near-above bin (panel 1 of 1)}
\begin{threeparttable}
\begin{tabular}{lrrrrrr}
\toprule
Date & N & Firms & Ref. N & Ref. firms & Coefficient & 95\% CI \\
\midrule
2016H2 & 43 & 43 & 33 & 33 & -0.0235 & [-0.0436, -0.0033] \\
2017H1 & 41 & 41 & 33 & 33 & -0.0230 & [-0.0382, -0.0077] \\
2017H2 & 44 & 44 & 33 & 33 & -0.0201 & [-0.0344, -0.0059] \\
2018H1 & 43 & 43 & 34 & 34 & -0.0171 & [-0.0304, -0.0038] \\
\bottomrule
\end{tabular}
\begin{tablenotes}[flushleft]\footnotesize
\item[] Outcome is cumulative MI change since end-2019 / end-2019 total assets, with firm and date fixed effects. Reference date: 2019H2; reference bin: $g < -0.10$. Firm-clustered SEs with t critical values (model firms minus one degrees of freedom); pointwise CIs, no multiple-testing correction. Complete results for four outcomes and both anchors, including normalizations, are in the CSV and figure. For the end-2019 bin $0 \leq g < 0.10$, pre-policy MI coefficients are 2016H2: -0.0235 [-0.0436, -0.0033]; 2017H1: -0.0230 [-0.0382, -0.0077]; 2017H2: -0.0201 [-0.0344, -0.0059]; 2018H1: -0.0171 [-0.0304, -0.0038]. These pointwise intervals exclude zero: firms just above the threshold were already building MI faster than the reference group before 2020. Part B describes differential adjustment, not a policy effect. Mechanical accounting comparisons are not attainable balance sheets, causal effects, or economic leverage; minority interests may be genuine equity.
\end{tablenotes}
\end{threeparttable}
\end{table}

